\subsection{Motivation: Wrong Inputs Defeat Correct Calculators}
\label{sec:intro_motivation}

Clinical decision-support is increasingly delivered by large language
models (LLMs) that call deterministic backend tools --- risk calculators,
dosing engines, eligibility checks --- via a structured tool-calling
interface. The deterministic backend is the easy part: a correctly
implemented CHA\textsubscript{2}DS\textsubscript{2}-VASc calculator
returns the right score on any input. The hard part is the boundary
where the LLM hands the calculator its inputs. A correct calculator
called with a transposed creatinine value, the wrong patient sex, or a
miscoded comorbidity will return a plausible-but-wrong answer that no
amount of downstream verification of the textual response can detect.

Concretely, in the medical\_calculator\_eval benchmark used in this study,
many vignettes admit only one ``correct'' calculator and one ``correct''
input mapping. A model that names the right calculator but misreads
the patient's sex from the case will receive a confidently wrong score
and present it as the answer. The failure mode is structural: it lives
at the LLM-to-tool input boundary, not in the calculator and not in
the final textual summary.

\subsection{Scope: Procedural Input Verification}
\label{sec:intro_scope}

We study \emph{procedural input verification} of tool calls: an
intermediate intervention that inspects the JSON payload the LLM sends
to a deterministic clinical calculator \emph{before} the calculator
runs. This is distinct from output verification (re-checking the final
textual answer), end-to-end outcome verification (was the diagnosis
right?), and retrieval grounding (was the cited fact in the source?).
The verification fires at one specific point in the agent loop: between
the model emitting a \verb|<tool_call>| block and the MCP tool being
invoked. The verifier compares each (field, value) pair the model
proposes to dispatch against an independently extracted reference set
of patient facts. If the case (and the extractor) say sex = female and
the model proposes sex = male, the call is flagged and the model is
asked to revise before the calculator runs.

This is the scope of every claim in this paper. We do not study other
verification surfaces, and our results say nothing about them directly.

\subsection{The Verification Hypothesis}
\label{sec:intro_hypothesis}

We study a single experimental factor: the architecture of the
input-verification pipeline. Four primary hypotheses are
pre-registered, one per architectural axis (\cref{sec:methods_stats}):

\textbf{H1 (placement).} Best-practice procedural input verification
with upfront full-schema extraction improves accuracy over the
forced-tool baseline B$'$.

\textbf{H2 (placement).} Best-practice procedural input verification
with reactive per-tool-call extraction improves accuracy over B$'$.

\textbf{H3 (output format).} Adding citation grounding to the
reactive pipeline (\texttt{value} + verbatim source span) improves
accuracy over the bare-value reactive baseline.

\textbf{H4 (sampling).} Replacing single-sample reactive extraction
with $k=3$ majority voting improves accuracy over the single-sample
reactive baseline.

``Best-practice'' here means schema-gated verifier comparison,
query-style intervention prompts, and prompt-only calibrated
abstention --- the three hygiene properties baked into all primary
study conditions (\cref{sec:methods_prereg}). These properties are
not factors under test; they were motivated by pilot diagnostics
(\cref{sec:methods_pilots}) and are held constant across H1--H4.

\subsection{Contributions and Roadmap}
\label{sec:intro_contributions}

This paper contributes the following. (i) A nine-condition evaluation
of procedural input verification on a clinical tool-calling benchmark
under a best-practice pre-registered design that varies only the
extractor pipeline across the four primary study conditions.
(ii) Six pre-registered McNemar contrasts (Bonferroni
$\alpha = 0.00833$) testing four hypotheses about extractor
architecture: placement (upfront vs reactive), output format (bare vs
citation-anchored), and sampling ($k=1$ vs $k=3$ majority vote).
(iii) A pilot-diagnostic methodology section that documents the
un-gated and correction-style variants that motivated the
best-practice hygiene baked into the primary study, including a
58.6\% non-required-field flag rate under un-gated extraction.
(iv) A uniform per-condition logging schema and analysis surface that
makes the four extractor pipelines directly comparable on accuracy,
flag composition, and cost.

We do \emph{not} claim that interwhen ``does not work.'' Our results
adapt the mechanism to a setting where the verifier's reference must be
an independent fact set extracted from the case --- a ground-truth
reference layer that interwhen's trace-state extraction does not have to
produce --- and the binding constraint, if one remains under
best-practice configuration, is jointly that reference pipeline and the
verifier's comparison surface.

The rest of the paper is organized as follows. \Cref{sec:bg_interwhen}
reviews the interwhen mechanism and the clinical-tool-calling setting.
\Cref{sec:methods_benchmark} describes the apparatus, the nine
conditions, the pre-registered configuration, and the statistical
analysis. \Cref{sec:results_gate} reports the apparatus gate and the
six pre-registered contrasts;
\cref{sec:results_architecture,sec:results_diagnostic} develop the
within-architecture comparisons and the post-hygiene flag
composition. \Cref{sec:disc_summary} synthesizes what the results
show and the implications for deploying intermediate verification in
clinical AI.
